\begin{textsample}{NEG Dim 8 – global\_north\_2025\_04 – Score -1.00 – global\_north\_2025\_04/124117329.txt}  \label{ex:f8_neg_419}
Palestinian medics mourn next to the bodies of colleagues killed on March 23, when Israeli forces opened fire on them as they were driving ambulances to help people wounded in an earlier Israeli attack in Rafah File : Hatem Khaled/Reuters A Palestinian paramedic who survived a deadly Israeli attack on a group of first responders in southern Gaza last month has been released from Israeli detention, the Palestine Red Crescent Society ( PRCS ) says.
Assaad al-Nassasra, an ambulance driver, was among at least 10 Palestinian detainees who were released into the Gaza Strip on Tuesday, the PRCS said.
The agency shared footage on social media that showed a visibly emotional al-Nassasra, dressed in a bright red PRCS jacket, embracing his colleagues after 37 days in Israeli detention.
His exact whereabouts had been unknown after the Israeli military opened fire on Palestinian first responders in the Rafah area of southern Gaza on March 23, killing 15 health workers in an attack that drew widespread outrage and calls for an independent investigation."@ @ @ @ @ @ @ @ @ @ the massacre of medical teams in the Tel Al-Sultan area of Rafah Governorate,"the PRCS said.
The first moments of colleague Asaad Al-Nsasrah ' s arrival and reunion with his teammates following his release today, after 37 days in detention by the occupation forces.
He had been arrested while performing his humanitarian duty during the massacre of medical teams in the Tel... **26 ; 1512 ; TOOLONG The PRCS reported last month that Israeli forces opened fire on the medics, who were driving in ambulances to assist wounded Palestinians at the site of an earlier Israeli attack.
The agency said it lost contact with its team and Israeli forces blocked access to the site of the incident.
When United Nations and Palestinian officials were able to reach the area a week later, they found a mass grave where bulldozed ambulances and bodies were buried.
Eight PRCS workers were killed along with six Palestinian Civil Defence team members and one UN employee, the PRCS said."This massacre of our team is a tragedy not only for @ @ @ @ @ @ @ @ @ @ for humanitarian work and humanity,"the agency said in a statement on March 30.
A video recovered from the mobile phone of one of the slain medics showed their final moments.
They were wearing highly reflective uniforms and were inside clearly identifiable rescue vehicles before they were shot by Israeli forces.
Amid the international outcry, the Israeli military announced it would investigate what happened.
It said last week that its probe had identified a series of"professional failures".
The army said its code of ethics was not violated and one soldier was dismissed.
The PRCS slammed the Israeli military ' s findings and called for an independent and impartial investigation by a UN body.
One of two survivors Al-Nassasra, 47, is one of two people who survived the attack.
The other survivor, Munther Abed, said at the time that he had seen al-Nassasra being captured, bound and taken away.
Advertisement The father of six last spoke to his family on the night of @ @ @ @ @ @ @ @ @ @ was on his way to the PRCS headquarters to break his Ramadan fast with his colleagues, according to his son Mohamed.
When the family tried to call him about dawn the next day, he didn't respond, and they found out from the PRCS that nobody could reach him or the other emergency workers.
Al-Nassasra had always warned his family that whenever he headed out on a mission, he may not make it back, his son said.
But the family tried not to think about that as al-Nassasra continued his work throughout Israel ' s 18-month war on Gaza.
His colleague Ibrahim Abu al-Kass also told Al Jazeera that al-Nassasra always carried sweets to offer children to encourage them to play somewhere safe, not in the middle of the road.
Israel has carried out an intensified campaign of arrests during the war.
According to the Palestinian prisoner support network Addameer, at least 9,900 Palestinians are currently being held in Israeli detention facilities, including 400 children.
More than 3,400 are held without @ @ @ @ @ @ @ @ @ @ detention", which can be renewed for \textbf{six-month} periods indefinitely.
Al-Nassasra was released into Gaza through the Kissufim checkpoint along with the 10 other detainees before they were sent to a hospital in central Gaza ' s Deir el-Balah for medical checkups.
Reporting from the city, Al Jazeera ' s Tareq Abu Azzoum said the released detainees reported being tortured in"horrific ways"and were in a bad physical and psychological state.
Advertisement Israeli forces have routinely targeted first responders, humanitarian workers and journalists during the Gaza bombardment.
More than 52,300 Palestinians have been killed since the war began on October 7, 2023, while at least 117,905 have been wounded, according to Gaza ' s Ministry of Health.

% matched lemmas: six-month
\end{textsample}
